% =====================================================================
%  9. ABSTRACT
%  Heading ALL CAPS 14pt bold; body 12pt, 1.5 spacing, indented
%  paragraphs, justified.
%
%  HARD LIMIT: the abstract MUST fit within two pages. The grading
%  rubric explicitly rewards "a good summary in two pages or less", so
%  after writing the real text, check the PDF and trim if it spills onto
%  a third page.
%
%  The \addcontentsline inside \fypfrontpage puts "ABSTRACT" into the
%  Table of Contents with its roman-numeral page number, as required.
% =====================================================================
\fypfrontpage{ABSTRACT}

% ---------------------------------------------------------------------
%  A results-focused abstract usually has five moves. Keep one short
%  paragraph per move and you will land comfortably inside two pages.
% ---------------------------------------------------------------------

% MOVE 1 -- CONTEXT
% TODO: two or three sentences establishing the domain the project sits
% in and why that domain matters.
[CONTEXT: what field does this project belong to, and why does it
matter? Write two or three sentences, no citations needed here.]

% MOVE 2 -- PROBLEM
% TODO: state the specific gap or limitation this project addresses.
[PROBLEM: the specific limitation of existing systems or practice that
this project set out to solve. This should be the same problem you
justify at length in Section 1.1 and back up with the literature in
Chapter 2.]

% MOVE 3 -- APPROACH
% TODO: name the method, architecture and main tools in one paragraph.
[APPROACH: the methodology followed, the architecture adopted, and the
principal technologies used. Name them concretely -- the reader should
be able to tell what was actually built.]

% MOVE 4 -- RESULTS  (the rubric weights this most heavily)
% TODO: give NUMBERS. Accuracy, response time, number of test cases
% passed, user-satisfaction score, throughput -- whatever you measured.
[RESULTS: the quantitative outcome. Do not write "the system works
well"; write what was measured and what the measured values were. An
abstract that is graded highly reports results, not intentions.]

% MOVE 5 -- CONCLUSION / SIGNIFICANCE
% TODO: one or two sentences on what the results mean and the single
% most promising direction for future work.
[SIGNIFICANCE: what the results demonstrate, and the one most valuable
avenue for follow-up work.]

\clearpage
